\chapter{Lösung}

\section{Hardware}

\subsection{Bauteilauswahl}

\subsubsection{Computer}

Für die Diplomarbeit sind 2 verschiedene Rechner in Frage gekommen. 

\begin{itemize}
\item \textbf{Raspberry Pi 5:} Der Raspberry pi 5 verfügt über einen Quad-Core ARM Cortex-A76 Prozessor wodurch genug Rechenleistung bereitgestellt ist. Zusätzlich bietet dieser genug Schnittstellen um mit einer einfach 8x8 Matrix die Schaltung zu entwickeln. Ein weiterer Punkt ist die gute Dokumentation, möglichen Librarys und die Programmiermöglichkeit mittels Pyhton welche für Diplomant Maximilian Koch bevorzugt ist.
Die Nachteile des Raspberry Pi sind zum einen den bnötigten Platz welchen dieser benötigt, die Anschaffungskosten für einen und der höhere Stromverbrauch.\\

\item \textbf{ESP32-S3:} Der ESP32-S3 verfügt über einen Dual-Core-Xtensa-LX7 Prozsessor welcher für das auslesen der Matrix genug Rechenleistung bereitstellen könnte. Für den ESP32-S3 spricht auch seine kompakte Größe, seine geringe Leisutngsaufnahme und die geringen Anschaffungskosten.
Die Nachteile des ESP32 sind zum einen aber auch die begrenzte Rechenleistung. Die weitere Hardware welche auch vom ESP gesteuert wreden muss, könnte den ESP schnell an seine Grenzen bringen und das fehlen eines vollwertigem Betriebssystem welches das Umsetzen von komplexeren Softwarestrukturen erschwert. 
\end {itemize}

Die Entscheidung fiel auf den Raspberry Pi 5, da dieser bereits verfügbar war und über ausreichende Rechenleistung  für das Projekt verfügt. Zudem stellte die Möglichkeit der Programmierung in Python, welche von Diplomant Maximilian Koch bevorzugt wird, einen weiteren ausschlaggebenden Faktor dar, wobei grundsätzlich auch andere Raspberry Pi Modelle für dieses Projekt geeignet gewesen wären.


\subsubsection{Sensoren}

\subsection{Schaltungsentwicklung}

\subsection{Platinendesign}

\subsection{Gehäusefertigung}

\subsection{Stromversorgung}


\section{Software}

\subsection{Magneterkennung}

\subsection{Matrixauslesung}

\subsection {Entwicklung eines Menüs}

\subsection{Spiellogik}

\subsubsection{Das Einzelspiel (Spieler gegen Computer)}

\subsubsection{Der Mehrspieler (Spieler gegen Spieler)}

\subsection{Einbindung des Spielererkennung}


\newpage

\section{Webinterface}

\subsection{Website}

\subsection{Datenbank}

\subsection{Server}



